% LaTeX preamble for "Dive into Deep Learning" PDF builds.
% Referenced by tools/gen_pdf.py via `include-in-header`.
% Quarto already loads fontspec/hyperref/amsmath/amssymb/graphicx/longtable/booktabs,
% so this file only adds project-specific tweaks.

% Allow line breaks inside long verbatim spans (URLs, code lines).
\usepackage{fvextra}
\DefineVerbatimEnvironment{Highlighting}{Verbatim}{commandchars=\\\{\},breaklines,breakanywhere,fontsize=\small}
% Highlighting (above) wraps source code, but pandoc emits cell *output* as the
% kernel's plain `verbatim`, which clips at the right margin. Redefine it as a
% breakable fvextra Verbatim so wide output — generated-text samples, dir() API
% dumps, profiler tables — wraps instead of running off the page. (Pandoc always
% puts \end{verbatim} on its own line, which fvextra requires.)
%
% This MUST be deferred to \AtBeginDocument. This file is injected via
% `include-in-header`, and Quarto appends its own preamble *after* it — that
% block loads tcolorbox (for callouts / the Shaded code box), which in turn
% \RequirePackage{verbatim}. verbatim.sty re-defines the `verbatim` environment
% unconditionally, silently reverting a definition made here at header time.
% That is why wide cell output kept overflowing the page edge (text lost, not
% wrapped) even with this line present: it was being overwritten before
% \begin{document}. \AtBeginDocument runs after every package, so ours wins.
\AtBeginDocument{%
  \DefineVerbatimEnvironment{verbatim}{Verbatim}%
    {breaklines,breakanywhere,fontsize=\small}%
  % Let pages end short rather than stretching the glue between code blocks,
  % output blocks and figures. book.cls leaves \flushbottom in force for
  % `twoside`, which on code-dense pages inflates every inter-block gap to
  % centimetres of white space just to square off the bottom margin. Facing
  % pages may now end at slightly different heights — the accepted trade in
  % technical books, and far less distracting than the stretch.
  \raggedbottom
}

% The book passed 1000 pages, and LaTeX's default page-number slot in the table
% of contents (\@pnumwidth = 1.55em) fits three digits. Every entry pointing at
% a four-digit page overflows its box by exactly 4.93pt, which accounted for 653
% of the build's Overfull \hbox warnings. Widen the slot, and the matching right
% margin with it (\@tocrmarg must stay about 1em clear of \@pnumwidth or the dot
% leaders run into the number).
%
% NOTE: \makeatletter must sit OUTSIDE any \AtBeginDocument wrapper. TeX fixes
% catcodes when it *reads* a token list, so a \makeatletter placed inside the
% braces comes too late: \@pnumwidth would already have been tokenized as \@
% followed by the letters "pnumwidth", silently redefining \@ instead. These
% assignments do not need deferring in the first place, since the TOC is
% typeset later, at \tableofcontents.
\makeatletter
\renewcommand{\@pnumwidth}{2.4em}%
\renewcommand{\@tocrmarg}{3.4em}%
\makeatother

% Better float placement for figures inside narrow sections.
\usepackage{float}

% mathtools (extends amsmath) provides the small parenthesized-matrix
% environments — pandoc's LaTeX writer normalizes
% \left(\begin{smallmatrix}…\end{smallmatrix}\right) into \begin{psmallmatrix},
% which is undefined without this. Load after amsmath (which Quarto loads).
\usepackage{mathtools}

% Allow underscores in identifiers without escaping.
\usepackage[strings]{underscore}

% Show subsections in the TOC.
\setcounter{tocdepth}{1}

% Looser hyphenation to reduce overfull \hbox warnings in code-heavy paragraphs.
\sloppy

% Chinese characters. The Language Models part prints Chinese in code cells
% and cell outputs, and the book's fonts (Source Serif 4, Source Sans 3,
% Inconsolata) have no CJK glyphs, so XeTeX would drop them ("Missing
% character"). Each run of CJK characters is set in a subset of Noto Sans SC
% (static/fonts/NotoSansSC-Regular-subset.otf, built by
% tools/make_cjk_font_subset.py; SIL Open Font License, see
% NotoSansSC-LICENSE.txt). The switch uses XeTeX's inter-character tokens
% rather than xeCJK, which TeX Live hosts may not have and which also
% changes punctuation and spacing elsewhere: the CJK code points get a class
% of their own, entering that class from any other class opens a group with
% the CJK font, and leaving it closes the group. A run that starts inside a
% brace group ends before the group closes (XeTeX inserts the closing tokens
% before the `}`), so the groups always nest. Text without CJK characters,
% that is all Latin text, is set exactly as before. Keep the code-point
% ranges in step with CJK_RANGES of tools/make_cjk_font_subset.py.
\usepackage{fontspec}
\newfontfamily\bookcjkfont{NotoSansSC-Regular-subset.otf}[Path=static/fonts/]
\newXeTeXintercharclass\bookcjkclass
\makeatletter
\def\book@cjkrange#1#2{%
  \count@=#1\relax
  \loop
    \XeTeXcharclass\count@=\bookcjkclass
  \ifnum\count@<#2\relax\advance\count@\@ne\repeat}
\book@cjkrange{"3000}{"303F}% CJK symbols and punctuation
\book@cjkrange{"3400}{"4DBF}% CJK unified ideographs extension A
\book@cjkrange{"4E00}{"9FFF}% CJK unified ideographs
\book@cjkrange{"F900}{"FAFF}% CJK compatibility ideographs
\book@cjkrange{"FF00}{"FFEF}% halfwidth and fullwidth forms
% Every other class: 0 (ordinary, all characters by default), 4095 (boundary:
% glue, kerns, `{`, `}`, the start and end of a paragraph) and any class a
% package allocates.
\count@=\z@
\loop
  \ifnum\count@=\bookcjkclass\else
    \XeTeXinterchartoks\count@\bookcjkclass={\begingroup\bookcjkfont}%
    \XeTeXinterchartoks\bookcjkclass\count@={\endgroup}%
  \fi
\ifnum\count@<4095 \advance\count@\@ne\repeat
\makeatother
\XeTeXinterchartokenstate=1
